\documentclass[11pt,a4paper]{article}
\usepackage[utf8]{inputenc}
\usepackage{amsmath,amssymb,amsthm}
\usepackage{mathrsfs}
\usepackage{geometry}
\geometry{margin=1in}
\usepackage{hyperref}
\usepackage{enumitem}

\newtheorem{definition}{Definition}[section]
\newtheorem{theorem}{Theorem}[section]
\newtheorem{proposition}{Proposition}[section]
\newtheorem{remark}{Remark}[section]
\newtheorem{lemma}{Lemma}[section]

\title{\textbf{Cohomology in Algebraic Quantum Field Theory}\\[6pt]
\large Lecture by John E.\ Roberts}
\author{Lecture Notes from the AQFT 50th Anniversary Conference}
\date{}

\begin{document}
\maketitle

\begin{abstract}
John Roberts surveys the role of cohomology in algebraic quantum field theory,
beginning with an analogy to the classification of locally trivial $G$-bundles
and developing the theory of net cohomology and its application to superselection
theory. He discusses the extension to curved spacetimes, where the topology of
the underlying manifold influences the superselection structure, and argues that
cohomology has become the preferred tool for tackling problems in superselection
theory.
\end{abstract}

\tableofcontents

%---------------------------------------------------------------------
\section{Introduction to Cohomology}
%---------------------------------------------------------------------

Cohomology is a pervasive tool in mathematics, appearing in essentially every
mathematical theory. Although it comes in many varieties, there are unifying
features. Roberts illustrates these with an example closely analogous to its
use in superselection theory.

\subsection{Locally Trivial $G$-Bundles}

Consider a locally trivial bundle $F$ over a base space $B$ with structure
group $G$. Local triviality means there exists an open covering $\{O_i\}$
of $B$ and isomorphisms $\phi_i: F|_{O_i} \xrightarrow{\sim} O_i \times G$.
Computing transition functions:
\begin{equation}
  \phi_i \circ \phi_j^{-1}(x, g) = (x,\, z_{ij}(x)\cdot g),
  \quad x \in O_i \cap O_j,
\end{equation}
one obtains a \textbf{1-cocycle} $z_{ij}$ satisfying:
\begin{equation}
  z_{ij}\, z_{jk} = z_{ik}
  \quad\text{whenever } O_i \cap O_j \cap O_k \neq \emptyset.
\end{equation}

Two 1-cocycles $z$ and $z'$ are \textbf{cohomologous} if there exist
$y_i \in G$ such that $z'_{ij}\, y_j = y_i\, z_{ij}$, corresponding
to equivalence of $G$-bundles.

\subsection{General Features of Cohomology}

From this example, we extract typical features:
\begin{enumerate}[nosep]
  \item Cohomology has a \textbf{degree} (here $n=1$).
  \item In low degree, cohomology \textbf{classifies} objects
        ($H^1(B,G)$ classifies locally trivial $G$-bundles).
  \item It is a cohomology of a \textbf{mathematical object} (first variable $B$)
        with \textbf{coefficients} (second variable $G$).
  \item It may provide examples, aid the general theory, or prove to be
        merely an alternative language.
\end{enumerate}

%---------------------------------------------------------------------
\section{Net Cohomology and Superselection Theory}
%---------------------------------------------------------------------

\subsection{The Superselection Criterion}

In AQFT, one considers representations $\pi$ of the observable net which
are equivalent to the vacuum representation $\pi_0$ when restricted to
the spacelike complement of any double cone $O$:
\begin{equation}
  \pi|_{\mathcal{O}'} \simeq \pi_0|_{\mathcal{O}'}
  \quad\text{for all double cones } \mathcal{O}.
\end{equation}
This means there exist unitaries $V_\mathcal{O}$ intertwining $\pi$ and $\pi_0$,
and setting $\rho(A) = V\pi(A)V^*$ yields a localized endomorphism of the net.

\subsection{Simplicial Structure}

Roberts introduces a simplicial structure on the set $K$ of double cones:
\begin{itemize}[nosep]
  \item A \textbf{0-simplex} is an element $a \in K$.
  \item A \textbf{1-simplex} $b$ consists of three elements $(d_0 b, d_1 b, |b|)$
        with $d_0 b, d_1 b \leq |b|$ (the support of $b$).
  \item A \textbf{2-simplex} $c$ consists of compatible 1-simplices forming a triangle.
\end{itemize}

Setting $z(b) = V(d_0 b)\, V(d_1 b)^*$, one obtains a \textbf{1-cocycle}:
\begin{equation}
  z(d_1 c) = z(d_0 c)\, z(d_2 c),
\end{equation}
with $z(b) \in \mathfrak{A}(|b|)$ by Haag duality. This is \emph{net cohomology}---a
cohomology with local coefficients, analogous to sheaf cohomology.

\subsection{Endomorphisms and Cocycles}

The localized endomorphisms $\rho$ and the 1-cocycles $z$ are related by:
\begin{equation}
  z(b) \;\text{intertwines}\; \rho(d_1 b) \to \rho(d_0 b).
\end{equation}
In the simplest cases (when $K$ is directed), the categories of localized
endomorphisms and of 1-cocycles are equivalent.

%---------------------------------------------------------------------
\section{Early Applications of Cohomology}
%---------------------------------------------------------------------

\subsection{Solitons in Two Dimensions (1976)}

In two spacetime dimensions, a 1-cocycle of the observable algebra $\mathfrak{A}$
is also a 1-cocycle of the field algebra $\mathfrak{F}$, but with two possible
localizations (left or right). If $\rho_L = \rho_R$ for cocycles of $\mathfrak{A}$,
this need not hold for cocycles of $\mathfrak{F}$---the discrepancy gives rise
to the \textbf{soliton phenomenon}.

\subsection{Completeness of Sectors (1980)}

For a free field with gauge group $G$, the question arises whether all
superselection sectors correspond to irreducible representations of $G$.
Roberts identified two cohomological conditions (A and B) sufficient to
guarantee completeness.

\subsection{Essential Duality}

Essential duality ($\mathfrak{A}'' = \mathfrak{A}^d$) implies that the
superselection sectors for $\mathfrak{A}$ are equivalent to those for
$\mathfrak{A}^{dd}$, considered as $C^*$-categories. Wedge duality can
imply essential duality.

%---------------------------------------------------------------------
\section{Superselection Theory on Curved Spacetimes}
%---------------------------------------------------------------------

\subsection{Topology and Superselection Structure}

On curved spacetimes, the topology and causal structure of the manifold
can influence the superselection structure. Key results (from work with
Brunetti, Guido, Longo, and Ruzzi):
\begin{itemize}[nosep]
  \item In $d \geq 3$, the causal complement of 1-simplices has the same
        number of components as in Minkowski space: \textbf{no new solitonic
        phenomena} in higher dimensions.
  \item The theory of sectors goes through when the set of regular diamonds
        $K$ is directed.
  \item The interesting (and open) case is when $K$ is \emph{not} directed.
\end{itemize}

\subsection{Homotopy and Path Independence}

Two notions of path arise: the topological one and the combinatorial one
(concatenation of 1-simplices). Under mild conditions, the two fundamental
groups coincide:

\begin{theorem}[Ruzzi]
If $M$ is arcwise connected and Hausdorff, and $K$ is a basis for the
topology consisting of arcwise and simply connected subsets, then
$\pi_1^{\mathrm{top}}(M) \cong \pi_1^K(M)$.
\end{theorem}

A 1-cocycle that is path-independent yields a unitary representation of
the fundamental group.

\subsection{The Punctured Spacetime Approach}

When Cauchy surfaces are compact, charges cannot be transported to infinity.
The solution: remove a point $x$ from the Cauchy surface, forming the
\emph{causal puncture} $K(x)$. The restriction to $K(x)$ behaves like
Minkowski space, enabling:
\begin{itemize}[nosep]
  \item Construction of endomorphisms $\rho_z(x)$ of $\mathfrak{A}^\perp(x)$.
  \item Definition of the tensor product structure.
  \item The full DHR machinery (conjugates, symmetry, statistics).
\end{itemize}

\begin{theorem}[Ruzzi]
The restriction functor $F(x): Z^1_t(\mathfrak{A}) \to Z^1_t(\mathfrak{A}(x))$
is full and faithful. It preserves irreducibility.
\end{theorem}

The open question: are there sectors on $M(x)$ that do not extend to $M$?

\subsection{Tensor Products Without Directed Sets}

A key advantage of the cohomological approach: the tensor product of
1-cocycles can be defined even when $K$ is not directed, requiring only
connectedness:
\begin{equation}
  (z \otimes z')(b) = z(b)\, \mathrm{Ad}_{z(p)}\bigl(z'(b)\bigr),
\end{equation}
where $p$ is a suitable path beginning spacelike to the support of $b$
and ending in $d_1(b)$.

%---------------------------------------------------------------------
\section{Conclusion}
%---------------------------------------------------------------------

Cohomology has evolved into the preferred tool for superselection theory:
\begin{itemize}[nosep]
  \item Localized endomorphisms are simpler when $K$ is directed, but
        require domains of definition when it is not.
  \item 1-cocycles avoid domain issues and allow direct definition of
        tensor products.
  \item On curved spacetimes, path-dependent cocycles arise from modified
        superselection criteria, producing sectors that depend on the
        topology of the underlying spacetime.
  \item Recent work by Brunetti and Ruzzi extends superselection theory
        to the locally covariant framework using cohomological methods.
\end{itemize}

\end{document}
